\begin{workedexample}[Worked Example: Inductance of a ferrite toroid]
A toroidal ferrite core has relative permeability $\mu_r = 125$, cross-sectional
area $A = 0.10\ \mathrm{cm^{2}} = 1.0\times10^{-5}\ \mathrm{m^{2}}$, and mean
magnetic path length $l = 3.0\ \mathrm{cm} = 0.030\ \mathrm{m}$. Wound with
$N = 20$ turns, its inductance is
\[ L = \frac{\mu_0\,\mu_r\,N^{2}A}{l}
     = \frac{(4\pi\times10^{-7})(125)(20^{2})(1.0\times10^{-5})}{0.030}. \]
Evaluating: $\mu_0\mu_r = 1.571\times10^{-4}$; times $N^2A = 4.0\times10^{-3}$
gives $6.28\times10^{-7}$; divided by $l$ gives
\[ L = 2.09\times10^{-5}\ \mathrm{H} = 20.9\ \mu\mathrm{H}. \]
Without the core ($\mu_r = 1$) the same winding would give only $0.17\ \mu\mathrm{H}$
--- the ferrite multiplies inductance by $\mu_r$.
\end{workedexample}

\begin{keytakeaways}
\begin{itemize}[leftmargin=1.4em,itemsep=2pt]
\item Permeability multiplies inductance ($L \propto \mu_r$) and, because
$L \propto N^{2}$, inductance scales with the \emph{square} of the turns.
\item Ferrites combine high $\mu_r$ with high resistivity, which suppresses the
eddy-current losses that make laminated iron unusable at RF.
\item Cores saturate: beyond the rated flux $\mu$ collapses. Ferrite beads
deliberately add frequency-dependent \emph{loss} to damp common-mode RF.
\end{itemize}
\end{keytakeaways}

\begin{exercises}
\item You double the number of turns on a fixed core. By what factor does the inductance change?
\item Why is laminated silicon-steel unsuitable for an inductor at $100\ \mathrm{MHz}$?
\item A ferrite bead is rated to present about $100\ \Omega$ at $100\ \mathrm{MHz}$.
At that frequency is the impedance mainly reactive or resistive, and why?
\end{exercises}

\furtherreading{Hayward~\cite{hayward} (ch.~2); Ott~\cite{ott} (ferrites for EMC).}
